\section{Properties of the Revenue Function}
\label{sec:revenue-props}

In our market setting, the sellers are in a pricing game,
where each seller's payoff is given by her revenue function $R_j$.
Analyzing this pricing game, therefore, requires understanding the revenue function, which is the topic of this section.

\subsection{A Closed Form for the Total Revenue}
\label{sec:tot-rev-closed-form}

Analyzing equilibria is often simpler for \emph{fixed-sum games},
i.e., where the sum of the players' payoffs does not depend on the players' actions.
For example, for two players, mixed (non-pure) Nash equilibria in \emph{fixed-sum games}
can be efficiently computed using linear programming \cite{dantzig1951proof}.

For our market model, the pricing game among the sellers is similar in spirit to a fixed-sum game,
since the sellers compete for the same buyers' budgets.
Indeed, if prices are high enough to exhaust all buyers' budgets,
the sellers' revenues would sum to $B$, the total buyer budget.
For lower prices, however, buyers may not exhaust their budgets,
making the game non-fixed-sum. We study the total revenue
$R_{\all}(\pvec, \xvec) \defeq \sum_{j=1}^m R_j(\pvec, \xvec)$
and give a closed-form expression for it,
which precisely quantifies this difference between the low-price and high-price regimes.

\begin{observation}
\label[observation]{thm:rev-vs-F}
For any prices $\pvec \in \mathbb{R}_{\ge 0}^m$,
\begin{tightenum}
\item (Non-rivalrous goods) If $\xvec \in \DemNR(\pvec)$,
    we have $R_{\all}(\pvec, \xvec) = F(\sum_{j=1}^m p_j)$,
    where $F(y) \defeq \sum_{i=1}^n \min(b_i, y) \le \min(B, ny)$.
\item (Rivalrous goods) If $\xvec \in \DemR(\pvec)$,
    we have $R_{\all}(\pvec, \xvec) = \min(B, \sum_{j=1}^m p_j)$.
\end{tightenum}
In either case, we get $R_{\all}(\pvec, \xvec) \le B$.
\end{observation}
\begin{proof}
In the non-rivalrous setting, each buyer $i \in [n]$ either exhausts her budget
or purchases every good completely, contributing $\min(b_i, \sum_{j=1}^m p_j)$
to the total revenue. Thus, $R_{\all}(\pvec, \xvec) = F(\sum_{j=1}^m p_j)$.

In the rivalrous setting, the buyers collectively either exhaust their budgets
or purchase every good completely
(since utilities are strictly monotone and $\xvec$ is Pareto optimal).
Hence, $R_{\all}(\pvec) = \min(B, \sum_{j=1}^m p_j)$.
\end{proof}

Since $R_{\all}(\pvec, \xvec)$ does not depend on $\xvec$ when $\xvec$ lies in the demand set,
we simply write $R_{\all}(\pvec)$ instead when it is clear from context
whether the goods are rivalrous or not.
We also conclude that $R_{\all}$ is monotonic in prices.

\begin{observation}
\label{thm:rall-mono}
Let $\pvec, \qvec \in \mathbb{R}_{\ge 0}^m$ be price vectors
such that $\sum_{j=1}^m p_j \le \sum_{j=1}^m q_j$.
Then $R_{\all}(\pvec) \le R_{\all}(\qvec)$
(for both the rivalrous and non-rivalrous settings).
\end{observation}

\subsection{Unifying the Rivalrous and Non-Rivalrous Settings}
\label{sec:unify}

We now slightly refine the analysis of \cref{sec:tot-rev-closed-form}.

When prices are low enough that each good is fully sold,
each seller $j$'s revenue is $p_j$ in the rivalrous setting
and $n \cdot p_j$ in the non-rivalrous setting.

Now consider the regime where prices are large.
In the rivalrous setting, when the prices sum to at least $B$,
we get $R_{\all}(\pvec) = \min(B, \sum_{j=1}^m p_j) = B$.
In the non-rivalrous setting, when prices sum to at least $b_{\min}$,
then every buyer contributes at least $b_{\min}$ to the total revenue,
giving us $R_{\all}(\pvec) \ge n \cdot b_{\min}$.

The above properties capture almost everything about the revenue functions that we need
in our existence results (\cref{sec:exist}) and the dynamics (\cref{sec:dynamics}).
Note how both the rivalrous and non-rivalrous settings behave similarly,
with the differences showing up only in some constants
(the factor next to $p_j$: $1$ vs $n$, and the price sum threshold: $B$ vs $b_{\min}$).
Indeed, our analyses for the existence results and the dynamics
are the same for both settings up to these constants.

To save ourselves the trouble of writing these results separately for each setting,
we introduce parameters $\nu$, $\beta$, and $\kappa$, and use them hereafter
to perform a unified analysis of both settings.

\begin{center}
\begin{tabular}{c|c|c|c}
\text{Market} & $\nu$ & $\beta$ & $\kappa \defeq B/(\nu\beta)$
\\ \hline Rivalrous & 1 & $B$ & 1
\\ Non-rivalrous & $n$ & $b_{\min}$ & $b_{\avg}/b_{\min}$
\end{tabular}
\end{center}

We now restate the above properties of the revenue functions using these parameters.
We write $\Dem(\pvec)$ instead of $\DemR(\pvec)$ or $\DemNR(\pvec)$ to refer to the demand set
in general without specifying whether the goods are rivalrous or non-rivalrous.

\begin{observation}[Properties of the revenue function]
\label{thm:accounting}
In both the rivalrous and non-rivalrous settings,
for any prices $\pvec$ and any $\xvec \in \Dem(\pvec)$, we have
\begin{enumerate}
\item \label{item:rev-props:ub}For each seller $j \in [m]$,
    we have $R_j(\pvec, \xvec) \le \nu p_j$, i.e., her revenue is
    at most what she would earn by selling her good fully.
\item \label{item:rev-props:low-prices}If $\sum_{j=1}^m p_j \le \beta$, then every good
    is fully sold, so each seller $j$'s revenue $R_j(\pvec, \xvec)$ is $\nu p_j$.
\item \label{item:rev-props:high-prices}$\sum_{j=1}^m p_j \ge \beta \implies R_{\all}(\pvec) \ge \nu \beta$.
\end{enumerate}
\end{observation}
